% !TeX root = ../papers/nc-p3.tex
\section{非被約な二十点での階数二排除と平方根の不分岐性}
\label{sec:nonreduced-opposite-root-separation}
\begin{theorem}\label{thm:nonreduced-opposite-root-separation}
四辺有限な標準有限Frobenius窓の一位置の配置 $X$ に対し、全階数平方根の行列 $N(x)$ の
一般階数が2になることはない。特に定理\ref{thm:double-point-opposite-reconstruction}の
反対根の全一意性は非被約な $X$ にも成立する。階数2以下の軌跡が有限で三次評価が同型な根そのものは
正規化平方根スキームの被約孤立点である。
\end{theorem}
\begin{proof}[一般階数二の二種類の因子分解]
$N$ の十成分は二次式全体を張る。一般階数が2なら
$\bigwedge^2N=hgg^{\mathsf t}$ を前節と同じ原始表示で書き、
$g$ の次数は1か2である。定数の場合は全像が一つの二平面に含まれる。
次数1なら動く像二平面のPl\"ucker座標 $g$ が線型であり、Klein二次曲面内の線型空間を与える。
全像が固定三次元空間内なら十成分の独立性に反する。
従って各像が同じ直線を含む。これを最初の座標に取ってその商の対称三行行列に移ると、
六成分は独立で一般階数1である。原始表示と次数によりこれは $G G^{\mathsf t}$ で
$G$ の三成分は独立な一次式である。残る行の三成分は $G$ に比例し、
その三成分の最大公約元は1ゆえ $bG^{\mathsf t}$（$b$ は一次式）となる。
角の二次式を $f$ と書く。二次小行列式には $(f-b^2)G_iG_j$ がある。
$q=f-b^2$ は一般階数2なので非零である。$X$ 上の $q$ の生成するイデアルは一点の
極大イデアルの二乗で消え、巡回加群なので台スキームの長さは高々4である。

しかし文字空間 $U$ の非零な二次式は長さ4以下の台を持たない。
実際辺の完全跡で $U^2=(V^2)^\perp$ なので、
$Y=\operatorname{Spec}(D/\operatorname{ann} q)$ には $V^2$ の制限を消す非零汎関数がある。
長さ4以下の任意の部分スキームでは $V^2$ は全ての関数を張る。
長さ4の場合も直線上にないので単元を含む一次制限の次元は3以上であり、
全環に達するまで生成部分空間の積の次元は厳密に増える。長さ3以下でも同じである。
二次評価自体が単射であるため $q$ が全 $X$ で消えることもなく、矛盾である。

$g$ の次数2なら $h$ は非零定数である。分数体上の階数二次形式の非退化二次元部の判別式は平方である。
実際非零な主二次小行列式が $h g_i^2$ となり、基礎体は代数閉である。
従って $y^{\mathsf t}N(x)y$ は $y$ に線型な二因子に分かれる。
$x$ についての係数は原始なのでGaussの補題で多項式因子を取れる。
一因子が $x$ で定数なら一次元の固定因子を持つ二次形式しか張れない。
従って両因子は双一次で、$N(x)$ は二つの線型像 $Ax,Bx$ の対称積である。
全二次式の基点がないことから両写像は可逆である。
座標変更後 $N(x)=x\mathbin{\odot}D_0x$ とし $D_0$ は非スカラーとする。
この行列の二次小行列式は $x\wedge D_0x$ の各座標どうしの積（非零定数倍）である。
よって $X$ は $D_0$ の固有直線スキーム $Z=V(x\wedge D_0x)$ の
イデアルを二乗したスキームに含まれる。
\end{proof}

\begin{proof}[固有直線スキームの二重化の長さ上界]
異なる固有値があれば、各一般化固有空間（ベクトル次元 $n$）の射影部分に沿って
その外の線型座標は $Z$ のイデアルに入る。局所的に固有値を $\lambda$ とする
固有ベクトル方程式を立てると他ブロックの $D_0-\lambda$ が可逆なためである。
従ってこの位置に支持される $X$ の部分の長さは、その射影部分との交叉長の
$1+(4-n)$ 倍以下である。外の $4-n$ 座標の積は全て零だからである。
$n=1,2,3$ に対する交叉長上界はそれぞれ $1,3,7$。分割
$3+1,2+2,2+1+1,1+1+1+1$ で合計はそれぞれ $18,18,17,16$ 以下となる。

一固有値ならそのスカラーを引きJordan型を取る。単一四鎖の場合は一つの固有点の
アフィン座標で固有直線が $(y_2-y_1^2,y_3-y_1^3,y_1^4)$ の完全交叉であり、
二重化の長さは16である。型 $(2,1,1)$ なら固有直線イデアルは $\ell I_o$、
ここで $\ell=0$ が核平面、$o$ が像点である。$X$ の関数環で $\ell^2D$ は $I_o^2$ で消えるため
長さ4以下で、$D/\ell^2 D$ の長さは平面切断長の二倍14以下。合計18を超えない。

型 $(2,2)$ では $D_0$ の冪零部を $(x,y)\mapsto(y,0)$（各二座標）とする。
台の固定直線のチャート $x_1=1$ で $a=x_2,b=y_1,c=y_2-a b$ と置けば
$Z$ のイデアルは $(c,b^2)$ である。他のチャートも同じ。
二乗での法線単項式は $1,b,b^2,b^3,c,bc$ の六つである。
直線のイデアルによるフィルトレーションで、各次数の商は直線切断の関数環上
その次数の単項式数以下で局所生成される。よって全長は直線切断長の六倍18以下である。
型 $(3,1)$ では作用を $(x,y,z,w)\mapsto(y,z,0,0)$ とする。
固定直線 $J=(y,z)$ のうち $(1,0,0,0)$ 以外では固有直線イデアルも $J$ である。
その点のチャート $x=1$ で $a=w,b=y,c=z-b^2$ とすれば $I_Z=(c,b^3,ab)$、$J=(b,c)$ である。
$J^2/I_Z^2$ は長さ8である。$b^2,b^3$ の係数は $a^2$ まで、
$b^4,b^5,bc,b^2c$ の係数は $a$ までで尽くされる。
従って $X$ の長さは $J^2$ 切断で高々 $3\cdot3$、残りは8以下である。
全ての非自明な型で20に満たず、一般階数2を排除した。
\end{proof}

\begin{proof}[接方向の排除]
$M$ が有限な階数降下を持つ反対根、$A=k[\epsilon]/(\epsilon^2)$ とし、固定された
$X$ 上でその正規化根を一次変形する。平方写像の階数10は保存されるため、基底を持ち上げて
十変数の対称行列 $N(z)$ を可逆な座標変更として取れる。
$X_A$ 上の階数1条件により全余因子は $I_{2X}\otimes A$ に入る。
前節のシジジー持上げは $A$ 上でもそのまま成り立ち、
$P_A=\operatorname{Sym}(S_2)\otimes A$ の上で $J_N\to J_M\otimes A$ を与える。
根基を除いた階数2以下の共通台の補では両イデアルが単位イデアルであり、この射は一つの関数による乗法である。
この関数は余次元3を越えて延びる。実際二重数の各係数について多項式環のHartogs延長を使えばよい。
次数零なので $A$ の定数で、特殊化は非零ゆえ単元である。
補での等号は $P_A$ の両係数での稠密性により全体での等号となる。
十生成元の持上げで中山の補題を使えば $J_N=J_M\otimes A$ である。
従って座標変更は一般対称行列の階数2以下のスキームの射影安定化群に入る。
その被約群は特異台のVeroneseを保つので $\mathrm{PGL}_4$ であり、逆の包含は合同変換による。
標数零の有限型群スキームは滑らかなので二重数上にも余分な方向はない。

合同変換を戻せばVeronese座標は固定される。$X_A$ 上の元の文字との相違は単元倍のみで、
二乗して規格化空間を保つ。その一次係数 $h$ は $hU_M^2\subseteq U_M^2$ を満たす。
ここで $U_M$ は問題の反対座標空間であり三次評価は同型である。
系\ref{cor:opposite-square-root-collision}の互いに素な二一次式の証明により $h$ はスカラーである。
よってGrassmann接空間は零となる。この最後の被約性の主張では、実際の反対位置と同じく
三次評価が同型の全階数平方根を取る。一般階数二の排除自体には相手の三次評価を使わない。
\end{proof}
